\chapter{Endpoint e API}

\section{Criterio di dettaglio}

Le API riportate in questa sezione corrispondono ai path effettivamente esposti dai router del repository corrente. Dove le note iniziali in \texttt{utilities} descrivevano endpoint piu' ampi o ancora ipotetici, qui si privilegia la situazione reale del codice.

\section{API Gateway}

Il gateway espone il punto di ingresso unico per la SPA e applica il controllo JWT prima dei servizi protetti. I path effettivi sono:

\begin{itemize}
\item \texttt{/health};
\item \texttt{/ws/notifications} per il proxy real-time delle notifiche;
\item \texttt{/bff/dashboard};
\item \texttt{/bff/configurator/:sessionId};
\item proxy REST verso \texttt{/auth}, \texttt{/catalog}, \texttt{/cad}, \texttt{/cart}, \texttt{/orders}, \texttt{/notifications}, \texttt{/chatbot} e \texttt{/reports}.
\end{itemize}

\section{authenticationService}

I path reali del servizio sono:

\begin{itemize}
\item \texttt{POST /auth/register};
\item \texttt{POST /auth/login};
\item \texttt{POST /auth/guest};
\item \texttt{GET /auth/me};
\item \texttt{GET /auth/sessions};
\item \texttt{POST /auth/logout};
\item \texttt{DELETE /auth/sessions/:sessionId}.
\end{itemize}

\section{catalogService}

Il catalogo reale espone:

\begin{itemize}
\item \texttt{GET /catalog/health};
\item \texttt{GET /catalog};
\item \texttt{GET /catalog/:id};
\item \texttt{POST /catalog};
\item \texttt{PUT /catalog/:id};
\item \texttt{DELETE /catalog/:id}.
\end{itemize}

La vista frontend usa questi path attraverso il client centralizzato, filtrando per ricerca, categoria e disponibilita'.

\section{cadService}

Il CAD attuale e' completamente REST-based e non corrisponde alle vecchie bozze con sessioni collaborative WebSocket. I path reali sono:

\begin{itemize}
\item \texttt{GET /cad/health};
\item \texttt{POST /cad/configurations};
\item \texttt{GET /cad/configurations};
\item \texttt{GET /cad/configurations/:id};
\item \texttt{GET /cad/configurations/:id/next-options};
\item \texttt{PATCH /cad/configurations/:id/environment};
\item \texttt{PATCH /cad/configurations/:id/category};
\item \texttt{PATCH /cad/configurations/:id/column-plan};
\item \texttt{PATCH /cad/configurations/:id/design};
\item \texttt{POST /cad/configurations/:id/finalize};
\item \texttt{POST /cad/configurations/:id/reorder}.
\end{itemize}

\section{cartService}

Il carrello reale espone:

\begin{itemize}
\item \texttt{GET /cart/health};
\item \texttt{GET /cart};
\item \texttt{PUT /cart/items/:sku};
\item \texttt{DELETE /cart/items/:sku};
\item \texttt{DELETE /cart};
\item \texttt{POST /cart/checkout}.
\end{itemize}

La versione corrente non usa le vecchie rotte di popolamento manuale o order requests separate: il checkout e' il punto di passaggio verso l'order service.

\section{notificationService}

I path REST reali sono:

\begin{itemize}
\item \texttt{GET /notifications/subscriptions};
\item \texttt{POST /notifications/subscriptions};
\item \texttt{GET /notifications/subscriptions/:subscriptionId};
\item \texttt{PATCH /notifications/subscriptions/:subscriptionId};
\item \texttt{DELETE /notifications/subscriptions/:subscriptionId};
\item \texttt{GET /notifications};
\item \texttt{GET /notifications/unread/count};
\item \texttt{PATCH /notifications/:notificationId/read}.
\end{itemize}

\section{chatbotService}

I path reali sono:

\begin{itemize}
\item \texttt{GET /chatbot/health};
\item \texttt{POST /chatbot/sessions};
\item \texttt{GET /chatbot/sessions/:sessionId};
\item \texttt{GET /chatbot/sessions/:sessionId/messages};
\item \texttt{PATCH /chatbot/sessions/:sessionId/close};
\item \texttt{POST /chatbot/sessions/:sessionId/messages}.
\end{itemize}

Il socket client del frontend usa il path \texttt{/api/chatbot/socket.io} e scambia eventi \texttt{chat:message:send}, \texttt{chat:message:new}, \texttt{chat:typing} e \texttt{chat:error}.

\section{orderService}

I path reali osservati nel router sono:

\begin{itemize}
\item \texttt{GET /orders/health};
\item \texttt{POST /orders};
\item \texttt{GET /orders};
\item \texttt{GET /orders/:orderId};
\item \texttt{PATCH /orders/:orderId/cancel};
\item \texttt{PATCH /orders/:orderId/status}.
\end{itemize}

\section{reportingService}

Il servizio di reporting espone un'API volutamente ridotta:

\begin{itemize}
\item \texttt{GET /reports/health};
\item \texttt{GET /reports/trends/orders}.
\end{itemize}

La UI admin usa questo endpoint per visualizzare il trend ordini in un intervallo temporale.

\section{Note di coerenza}

Le note in \texttt{utilities} contenevano una mappa endpoint piu' ampia di quella realmente implementata per alcune aree. In questo documento sono stati riportati i path effettivi del repository corrente, evitando di descrivere rotte che non esistono nel codice. Questa scelta rende il documento direttamente verificabile.
